\section{Experimental Setup}
\label{results}

\subsection{Verification Tool}

We use Verus~\cite{Verus}, a semi-automated Rust verification tool with multiple prior applications at the intersection
of AI and verification \cite{AutoVerus,alphaverus,verussage,RAGVerus,vericode}. 

In Verus, users can write three kinds of code: 
\begin{itemize}
    \item \texttt{spec} code describes higher-level program properties such as function pre- and post-conditions
    or loop invariants.
    \item \texttt{proof} annotations give Verus hints for verifying proof goals:
    Verus is semi-automated, meaning it automatically proves many specifications, 
    but more complex proofs may require some help to determine intermediate proof goals 
    or useful lemma calls.
    \item \texttt{exec} code is Rust executable code. 
    During compilation, Verus erases spec and proof code, 
    leaving only the exec code: 
    thus, adding Verus proofs incurs no runtime performance cost.
\end{itemize}
Verus is designed to integrate directly into Rust code: 
it relies on the Rust type checker and uses Rust macros and annotations.
However, using Verus does involve some Verus-specific syntax, 
such as additional mathematical keywords or named function outputs to describe their expected properties.

\subsection{Experimental Design}
We limit this project's scope to reward hacking with few-shot prompts,
since this provides a preliminary comparison; 
we believe future RL training will amplify the results shown below.

\begin{enumerate}
    \item Our benchmark contains coding problems 
    with NL descriptions and formal Verus specs;
    we consider the formal Verus specs to be the ``gold-standard''
    ultimately defining the problems.
    \item We provide a model with only the NL descriptions, 
    and ask it to generate 
    (A) unit tests for evaluating this coding task and 
    (B) a formal Verus specification. 
    This generated spec may differ from the original gold-standard spec.
    \item For both (A) and (B), we check if the code produced also satisfies
    the gold-standard spec. 
    If it does not (and satisfies the unit tests/generated Verus spec), 
    we interpret this to mean the model successfully ``hacked''
    the reward mechanism. 
    We hypothesize we will see more reward hacking with unit tests,
    compared to the Verus generated specs.
\end{enumerate}
\begin{figure}[h]
\centering
\tikzset{
  box/.style={rectangle, rounded corners=3pt, draw, align=center,
              minimum width=2.0cm, minimum height=0.65cm, font=\small},
  goldbox/.style={box, fill=yellow!50, draw=orange!90!black, very thick},
  genbox/.style={box, fill=blue!30, draw=blue!80!black, very thick},
  rlbox/.style={box, fill=green!20, draw=green!60!black, very thick},
  evalbox/.style={box, fill=red!30, draw=red!80!black, very thick},
  arr/.style={-{Stealth[length=4pt]}, thick},
}
\begin{tikzpicture}[node distance=0.3cm and 1.0cm]

  % Top: benchmark input
  \node[goldbox, minimum width=3.4cm] (bench)
    {Benchmark\\{\footnotesize NL spec + gold Verus spec}};

  % NL spec fed to generator
  \node[rlbox, below=0.45cm of bench] (gen)
    {LLM Generator};
  \draw[arr] (bench) -- node[right, font=\footnotesize]{NL spec} (gen);

  % Two branches
  \node[genbox, below left=-.3cm and 1.2cm of gen]  (tests)
    {(A) Unit Tests};
  \node[genbox, below right=-.3cm and 1.2cm of gen] (spec)
    {(B) Generated Verus Spec};
  \draw[arr] (gen) -- (tests);
  \draw[arr] (gen) -- (spec);

  % RL training nodes
  \node[rlbox, below=0.35cm of tests] (rlA)
    {LLM Implementor\\{\footnotesize goal: pass unit tests}};
  \node[rlbox, below=0.35cm of spec]  (rlB)
    {LLM Implementor\\{\footnotesize goal: satisfy Verus spec}};
  \draw[arr] (tests) -- (rlA);
  \draw[arr] (spec)  -- (rlB);

  % Code outputs
  \node[box, below=0.35cm of rlA] (codeA) {Code A};
  \node[box, below=0.35cm of rlB] (codeB) {Code B};
  \draw[arr] (rlA) -- (codeA);
  \draw[arr] (rlB) -- (codeB);

  % Evaluation node
  \node[evalbox, minimum width=3.4cm,
        below=0.6cm of $(codeA)!0.5!(codeB)$] (eval)
    {Evaluate \\{\footnotesize reward hacking $\Leftrightarrow$ gold Verus spec not satisfied}};
  \draw[arr] (codeA) -- (eval);
  \draw[arr] (codeB) -- (eval);

  % Gold spec dashed arrow going around the right side
  \draw[arr, dashed, draw=orange!80!red!90!black, very thick] (bench.east)
        -- ++(4.2,0)
        |- (eval.east)
        node[pos=0.25, right, font=\small\bfseries, text=orange!80!red!90!black]{gold spec};

\end{tikzpicture}
\caption{Experimental setup. Both branches receive only the natural language specification as input. The dashed arrow indicates the gold-standard Verus specification used solely for final evaluation.}
\label{fig:setup}

\end{figure}

% Note: it's possible that the code produced by the models does satisfy the gold-standard specification,
%     but Verus rejects it due to insufficient proof annotations. 
%     If we find this to be the case, 
%     we would add a step between (3) and (4) which asks a model to add only proof annotations to the produced code,
%     to help Verus see that the implementation does satisfy the specification.

\textbf{Benchmark}: We use the HumanEval-Verus
dataset~\cite{alphaverus}, which has 78 coding tasks with
NL descriptions, formal Verus specs, and Rust
implementations satisfying the Verus specs. 
This is the largest known
benchmark with NL and formal Verus specs. 
We leave other benchmarks to future work.